% TOP_PROBLEM: {"problemId": "problem.alperins-weight-conjecture-for-modular-representations-of-finite-groups", "canonicalTitle": "Alperins Weight Conjecture for Modular Representations of Finite Groups", "releaseRank": 342, "primaryDomain": "algebra_representation_category", "primaryDomainLabel": "Algebra, representation & category theory", "displayStatus": "open_with_solved_subcases", "releaseStatus": "open_with_solved_subcases"}
% !TeX program = lualatex
% Definition and sources only; this file is not an LLM solution attempt.
\documentclass[11pt]{article}
\usepackage[margin=25mm]{geometry}
\usepackage{fontspec}
\setmainfont{Latin Modern Roman}
\usepackage{amsmath,amssymb,mathtools}
\usepackage{unicode-math}
\setmathfont{Latin Modern Math}
\usepackage{xurl,hyperref}
\hypersetup{colorlinks=true,urlcolor=blue}
\setcounter{secnumdepth}{0}
\setlength{\parindent}{0pt}
\setlength{\parskip}{5pt}
\setlength{\emergencystretch}{3em}
\begin{document}
\section*{\texorpdfstring{Alperins Weight Conjecture for Modular Representations of Finite Groups}{Alperins Weight Conjecture for Modular Representations of Finite Groups}}\label{342}
\subsection{Definitions and mathematical statement}
Let $B$ be a $p$-block of a finite group $G$ over a splitting field of characteristic $p$, and let $\ell(B)$ be the number of isomorphism classes of its simple modules. A weight is a pair $(Q,\varphi)$ where $Q$ is a $p$-subgroup and $\varphi$ is an irreducible ordinary character of $N_G(Q)/Q$ of $p$-defect zero, meaning its degree has the full $p$-part of the group order. It is a $B$-weight when the block containing the inflated character corresponds to $B$ under the block-induction convention in the cited source. The assertion is
\[
 \ell(B)=\#\{G\text{-conjugacy classes of }B\text{-weights}\}.
\]
The local block-association condition is essential; counting all weights of $G$ instead would discard the selected blockwise statement. The conjecture asks for equality of finite numbers, not a specified canonical bijection.
\par\smallskip\textit{Formulation and concept references:} \href{https://par.nsf.gov/servlets/purl/10275354}{[S1]}.

\subsection{Short English statement}
Does each block have exactly as many simple modular representations as conjugacy classes of its associated local weights?
\subsection{Sources}
\begin{itemize}
\item[S1]Radical subgroups and the inductive blockwise Alperin weight conditions for PSp4(q), Introduction, printed page 1181, lines 20-24.
\url{https://par.nsf.gov/servlets/purl/10275354}.
\textit{Formulation source.}
\end{itemize}
\end{document}
